\documentclass[11pt]{article}
\usepackage[margin=1in]{geometry}
\usepackage{amsmath,amssymb,amsthm,booktabs,hyperref}
\hypersetup{colorlinks=true,linkcolor=blue,urlcolor=blue}
\newcommand{\Q}{\mathbb Q}
\newcommand{\sech}{\operatorname{sech}}
\newcommand{\Res}{\operatorname{Res}}
\newtheorem{theorem}{Theorem}
\title{All fifth-moment Berndt integrals:\\a corrected arithmeticity theorem\\MathDB 361365}
\author{Research note}
\date{17 September 2026}
\begin{document}
\maketitle

\section*{Task 1. Resolve the arithmeticity conjecture}
For each integer $M\ge1$, put
\[
 J_M=\int_0^\infty\frac{x^5\,dx}{(\cos x+\cosh x)^M},\qquad
 G=\Gamma(1/4),\qquad u=\frac{G^4}{4\pi^3}.
\]
We prove the even-power assertion of Xu--Zhao's Conjecture 5.9
\cite[p.~27]{XZ} for every $M=2n$. The odd-power assertion as printed
is false: its bulk upper limit $n-3$ omits an entire band. We prove
the corrected assertion, with upper limit $n-2$, for every $M=2n-1$,
and give an exact counterexample to the printed assertion at $M=7$.
Thus the result includes a full all-power theorem, beyond identifying
the indexing error. Empty sums below mean zero.

\begin{theorem}\label{support}
For every $M\ge1$, the number $J_M$ is a polynomial over $\Q$ in
$u$ and $\pi$. Every nonconstant monomial $u^j\pi^a$ occurring in
this polynomial satisfies
\begin{equation}\label{bounds}
 j>0,\qquad j\equiv M\pmod2,\qquad 1\le a\le6,
 \qquad j-a\le M-4.
\end{equation}
In addition, $j=1$ permits only $a=1,2$; $j=2$ permits only
$a=1,2,3,4$; and $(j,a)=(4,5)$ never occurs. A rational constant
is permitted only when $M$ is even.
\end{theorem}

In the notation of the source, this proves, for every $n\ge1$,
\begin{align}
 J_{2n-1}\in{}&\frac{G^4}{\pi}\Q+\frac{G^4}{\pi^2}\Q
 +\sum_{j=2}^{n-2}\sum_{i=1}^{6}
       \frac{G^{8j-4}}{\pi^{6j-10+i}}\Q\notag\\
 &+\sum_{j=1}^{3}\sum_{i=1}^{7-2j}
       \frac{G^{8j+8n-20}}{\pi^{6j+6n-22+i}}\Q,\label{odd-result}\\
 J_{2n}\in{}&\Q+\sum_{i=1}^{4}\frac{G^8}{\pi^{i+1}}\Q
 +\sum_{\substack{1\le i\le6\\i\ne2}}
       \frac{G^{16}}{\pi^{i+5}}\Q\notag\\
 &+\sum_{j=3}^{n-2}\sum_{i=1}^{6}
       \frac{G^{8j}}{\pi^{6j-7+i}}\Q
 +\sum_{j=1}^{3}\sum_{i=1}^{7-2j}
       \frac{G^{8j+8n-16}}{\pi^{6j+6n-19+i}}\Q.\label{even-result}
\end{align}
The bounds are inclusions, not assertions that every allowed
coefficient is nonzero. This distinction also handles overlaps in
the displayed spaces for small $n$.

\subsection*{1. Residues for an arbitrary denominator power}
Let $a_0=(1+i)/2$, $b_0=(1-i)/2$, and
$F(z)=z^5/(\cos z+\cosh z)^M$.
The factorization
\[
 \cos z+\cosh z=2\cosh(a_0z)\cosh(b_0z)
\]
places the first-quadrant poles at $z_n=(1+i)t_n$, where
$t_n=\pi(n+1/2)$ and $n\ge0$. Integrate on the positive triangle
with vertices $0,R,iR$, $R=2\pi N$. On its oblique side $x+y=R$,
\[
 |\cosh(a_0z)|=\cosh((x-y)/2)\ge1,\qquad
 |\cosh(b_0z)|\ge\sinh(R/2).
\]
That side contributes $O_M(R^6e^{-MR/2})$. Since $F(iy)=iF(y)$,
the two axial sides together contribute $2\int_0^R F(x)\,dx$.
Consequently
\begin{equation}\label{res-sum}
 J_M=\pi i\sum_{n\ge0}\Res_{z=z_n}F(z).
\end{equation}
The integral is convergent; the residue series obtained next is
absolutely convergent by exponential decay of hyperbolic secants
and their derivatives.

Define rational coefficients
\[
 A_{M,h}=[v^{2h}]\left(\frac{v}{\sinh v}\right)^M.
\]
Expanding $F(z_n+w)$ gives
\begin{equation}\label{res-formula}
 \begin{aligned}
 J_M={}&(-1)^M2^{3-M}\pi
 \sum_{h=0}^{\lfloor(M-1)/2\rfloor}(-1)^h A_{M,h}
 \sum_{\ell=0}^{\min(5,M-1-2h)}\frac{\binom5\ell}{r!}\\
 &\hspace{18mm}\cdot\sum_{n\ge0}(-1)^{nM}t_n^{5-\ell}
                (\sech^M)^{(r)}(t_n),
 \qquad r=M-1-2h-\ell.
 \end{aligned}
\end{equation}
Here and below a superscript in parentheses denotes ordinary
differentiation of the whole expression in parentheses. To check the phase in
\eqref{res-formula}, use
\[
 \cosh(a_0(z_n+w))=i(-1)^n\sinh(a_0w),\quad
 \cosh(b_0(z_n+w))=\cosh(t_n+b_0w),
\]
and, when $2h+\ell+r=M-1$,
\[
 2^{-M}i^{-M}a_0^{2h-M}(1+i)^{5-\ell}b_0^r
       =i(-1)^{M+1+h}2^{3-M}.
\]

Put $\delta=0$ for odd $M$ and $\delta=1$ for even $M$.
The identity
\[
 (\sech^m)''=m^2\sech^m-m(m+1)\sech^{m+2}
\]
implies
\begin{equation}\label{hyperbolic}
 \sech^M t=\sum_{k=0}^{(M-1-\delta)/2}C_{M,k}
       \left(\frac{d}{dt}\right)^{2k}\sech^{1+\delta}t,
 \quad
 \sum_k C_{M,k}T^{2k}
 =\frac{\prod_{\substack{1+\delta\le j\le M-2\\j\equiv1+\delta\ (2)}}
                    (j^2-T^2)}{(M-1)!}.
\end{equation}
An empty product is one.

\subsection*{2. Divisor interchange bounds the number of derivatives}
For nonnegative integers $P,Q$ of opposite parity set
\[
 \Phi_{P,Q}(q)=\sum_{k,v\ge1}k^Pv^Qq^{kv},\quad
 D=q\frac d{dq},\quad
 s=\min(P,Q),\quad \kappa=\frac{|P-Q|+1}{2}.
\]
All the sums converge absolutely for $|q|<1$. Interchanging $k,v$
if necessary yields
\begin{equation}\label{divisor}
 \Phi_{P,Q}=D^sL_{2\kappa-1},\qquad
 L_{2\kappa-1}(q)=\sum_{k\ge1}\frac{k^{2\kappa-1}q^k}{1-q^k}
   =\frac{B_{2\kappa}}{4\kappa}(1-E_{2\kappa}(q)).
\end{equation}
The Eisenstein normalization is $q=e^{2\pi i\tau}$.
For $\kappa\ge2$, the classical modular-form structure theorem
gives $E_{2\kappa}\in\Q[E_4,E_6]$, homogeneous of weight $2\kappa$
\cite[Proposition 4]{Z}. Rationality follows from the rational
Fourier coefficients and an invertible rational coefficient matrix
for the monomial basis. For $\kappa=1$ we use $E_2$ itself.
Ramanujan's identities \cite[Section 7]{R} give
\begin{equation}\label{D}
 DE_2=\frac{E_2^2-E_4}{12},\quad
 DE_4=\frac{E_2E_4-E_6}{3},\quad
 DE_6=\frac{E_2E_6-E_4^2}{2}.
\end{equation}
Thus the nonconstant part of \eqref{divisor} is a rational
polynomial of weight $2w$, $w=\kappa+s=(P+Q+1)/2$.
Its degree in $E_2$ is at most $s$ when $\kappa\ge2$, and at most
$s+1$ when $\kappa=1$.

For even $M$ define
\begin{equation}\label{H}
 H_{P,Q}=\Phi_{P,Q}(q)
  -(2^P+2^{Q+1})\Phi_{P,Q}(q^2)
  +2^{P+Q+1}\Phi_{P,Q}(q^4),\qquad q=e^{-\pi}.
\end{equation}
Selecting odd $k$ and giving $v$ the sign $(-1)^{v-1}$ shows that
this is $\sum_{k\text{ odd},v\ge1}k^Pv^Q(-1)^{v-1}q^{kv}$.
For odd $M$ define
\begin{equation}\label{T}
 T_{P,Q}=\frac{\Phi_{P,Q}(iq)-\Phi_{P,Q}(-iq)}{2i},
       \qquad q=e^{-\pi/2}.
\end{equation}
The coefficient multiplier is $\chi_4(kv)=\chi_4(k)\chi_4(v)$,
where $\chi_4$ is the nonprincipal character modulo four.
Expanding $\sech t=2\sum_{j\ge0}(-1)^je^{-(2j+1)t}$ and
$\sech^2t=4\sum_{v\ge1}(-1)^{v-1}ve^{-2vt}$ now gives, for $d\ge0$,
\begin{align}
 \sum_{n\ge0}t_n^P(\sech^2)^{(d)}(t_n)
    &=\pi^P2^{2+d-P}(-1)^d H_{P,d+1},\label{even-sum}\\
 \sum_{n\ge0}(-1)^nt_n^P\sech^{(d)}(t_n)
    &=\pi^P2^{1-P}(-1)^d T_{P,d}.\label{odd-sum}
\end{align}
All differentiations and interchanges here are justified by locally
uniform absolute exponential convergence for $t_n\ge\pi/2$.
In \eqref{res-formula}--\eqref{hyperbolic} these formulas use
\begin{equation}\label{indices}
 P=5-\ell,\quad d=M-1-2h-\ell+2k,\quad Q=d+\delta.
\end{equation}
In particular $0\le P\le5$, $Q\ge0$, $P,Q$ have opposite parity,
and $Q\ge1$ in the even case. Although $d$ grows with $M$, the
number $s$ of modular derivatives in \eqref{divisor} is at most five.

\subsection*{3. Exact modular substitutions and parity}
The classical lemniscatic value is $\theta_3(i)^4=u$
\cite[equation (2.3)]{XZ}. The required substitutions are
\begin{equation}\label{values}
\begin{array}{c|ccc}
 \tau&E_2&E_4&E_6\\\hline
 i/2&6/\pi-3u/2&33u^2/4&-189u^3/8\\
 i&3/\pi&3u^2/4&0\\
 2i&3/(2\pi)+3u/8&33u^2/64&189u^3/512\\
 (1+i)/4&12/\pi-3iu&-33u^2&189iu^3
\end{array}
\end{equation}
For completeness, theta duplication \cite{DLMF}, starting from
$\theta_2(i)^4=\theta_4(i)^4=u/2$, gives at $2i$
\[
 (\theta_2^4,\theta_3^4,\theta_4^4)
 =u\left(\frac{3-2\sqrt2}{8},\frac{3+2\sqrt2}{8},
          \frac1{\sqrt2}\right).
\]
Substitution in
\[
 E_4=\tfrac12(\theta_2^8+\theta_3^8+\theta_4^8),\quad
 E_6=\tfrac12(\theta_2^4+\theta_3^4)
              (\theta_3^4+\theta_4^4)(\theta_4^4-\theta_2^4)
\]
gives the last two modular weights at $2i$ and $i$. For weight two use
\[
 2E_2(2\tau)-E_2(\tau)=\tfrac12(\theta_3(\tau)^4+
                                      \theta_4(\tau)^4),\quad
 E_2(-1/\tau)=\tau^2E_2(\tau)+\frac{6\tau}{\pi i}.
\]
The latter gives $E_2(i)=3/\pi$. The modular transformations of
$E_4,E_6$ give the row $i/2$. Finally,
$(1+i)/4=-1/(2i-2)$ gives the last row by translation and inversion.
The values at $(-1+i)/4=-1/(2i+2)$ are obtained from this last row
by replacing $u$ with $-u$.

For a homogeneous polynomial of weight $2w$, its value in the first
row of \eqref{values} is $4^w$ times its value in the third row
with $u$ replaced by $-u$. Because $2^{P+Q+1}=4^w$, the first
and third terms of \eqref{H} have an even sum in $u$; the middle
term is also even. The constant part of $L$ when $s=0$ cancels:
in the even case $Q\ge1$, so $s=0$ forces $P=0$, and
$1-2^P-2^{Q+1}+2^{P+Q+1}=0$.
Thus $H_{P,Q}$ contains only even powers of $u$.
Equation \eqref{T} takes the odd part in $u$ of the last row,
divided by $i$, and therefore contains only odd powers of $u$,
with rational coefficients in $u,\pi^{-1}$.

\subsection*{4. Proof of every support restriction}
Consider a homogeneous modular monomial after substitution. Let $t$
be the number of selected $\pi^{-1}$ factors from the $E_2$ factors.
Its power of $u$ is $j=w-t$. Multiplication by the external factor
$\pi^{P+1}$ in \eqref{res-formula}, \eqref{even-sum}, and
\eqref{odd-sum} gives the power $a=P+1-t$. Consequently
\begin{equation}\label{ja}
 j=\frac{M+5+\delta}{2}-h-\ell+k-t,\quad
 a=6-\ell-t,\quad
 j-a=\frac{M-7+\delta}{2}-h+k\le M-4.
\end{equation}
For $\kappa\ge2$, $t\le s\le P$ implies $1\le a\le6$ and
$j\ge\kappa$. For $\kappa=1$, a positive $j=s+1-t$ likewise
forces $t\le s$. Parity was proved above.

The only possible $j=0$ terms have $\kappa=1$ and $t=s+1$.
They vanish in the odd case. In the even case their multiplier,
apart from a common rational factor and $\pi^{-s-1}$, is
\[
 3^{s+1}\left(2^{s+1}-2^P-2^{Q+1}+2^{P+Q-s}\right).
\]
When $P=Q+1$ this is zero. When $Q=P+1$ the external factor
$\pi^{P+1}$ makes $a=0$. This proves the assertion about constants.

If $j=1$, only $\kappa=1$ is possible. Then $t=s$, and
$|P-Q|=1$ gives $a=P+1-s\in\{1,2\}$.
If $j=2$, then either $\kappa=1$, $t=s-1$, giving
$a\in\{2,3\}$, or $\kappa=2$, $t=s$, giving $a\in\{1,4\}$.

It remains to exclude $j=4,a=5$. The equalities imply
\[
 a=P-s+5-\kappa.
\]
If $Q>P$, then $P=s$ and $a\le4$. Otherwise
$P-s=2\kappa-1$, so $a=5$ forces $\kappa=1$,
$P=Q+1$, and $t=s-3\ge0$. Since $P\le5$ and $Q\ge1$,
only $(P,Q)=(4,3),(5,4)$ remain. Direct use of
\eqref{divisor}, \eqref{D}, and \eqref{values} gives the exact identities
\begin{equation}\label{cancellations}
 H_{4,3}=\frac{3u^2}{16\pi^2},\qquad
 H_{5,4}=\frac{5u^2}{8\pi^3}.
\end{equation}
For a short way to reproduce this finite algebra, apply the
derivation in \eqref{D} three or four times to $(1-E_2)/24$,
then take the three substitutions with multipliers respectively
$(1,-32,256)$ or $(1,-64,1024)$. Neither result has a $u^4$ term.
This establishes all assertions of Theorem~\ref{support}.

To obtain \eqref{odd-result}, write $j=2b-1$. For $b\ge2$,
\eqref{bounds} permits all six $a$ values when $b\le n-2$;
for $b=n-1,n,n+1$ it requires $a\ge2,4,6$, respectively,
and no larger $b$ occurs. The $j=1$ restriction supplies the first
two terms. The same calculation with $j=2b$ gives
\eqref{even-result}, using the exceptional restrictions at $j=2,4$.
Any overlap for small $n$ enlarges the displayed space and causes
no problem. Finally use $u^j\pi^a=4^{-j}G^{4j}\pi^{a-3j}$.

\subsection*{5. An exact counterexample to the printed odd assertion}
The preceding formulas also give an effective exact evaluation for
each $M$. For $M=7$, they give
\begin{align}
 J_7={}&\frac{713\pi^6}{707788800}u^9
 +u^7\left(\frac{707\pi^6}{265420800}
        -\frac{7\pi^5}{442368}+\frac{103\pi^4}{5898240}\right)\notag\\
 &+u^5\left(\frac{51331\pi^6}{2123366400}
       -\frac{839\pi^5}{10616832}+\frac{581\pi^4}{2949120}
       -\frac{91\pi^3}{442368}+\frac{11\pi^2}{196608}\right)\notag\\
 &+u^3\left(\frac{75\pi^6}{524288}-\frac{1295\pi^5}{1179648}
       +\frac{114331\pi^4}{53084160}-\frac{1019\pi^3}{442368}
       +\frac{1253\pi^2}{884736}-\frac{7\pi}{24576}\right)\notag\\
 &+u\left(\frac{51331\pi^2}{8847360}
                         -\frac{839\pi}{221184}\right).\label{counter}
\end{align}
This calculation requires only $E_2,E_4,E_6,E_8=E_4^2$ in
\eqref{divisor}. The accompanying rational-arithmetic certificate
records all contributions indexed by $(h,k)$ in
\eqref{res-formula}--\eqref{hyperbolic}; their sum is exactly
\eqref{counter}. In particular the coefficient of $u^3\pi^6$ is
$75/524288\ne0$.

For $n=4$ the printed bulk sum $\sum_{j=2}^{n-3}$ is empty.
Its remaining terms involve only $u^1,u^5,u^7,u^9$, with rational
coefficients times powers of $\pi$. The established algebraic
independence of $\pi$ and $\Gamma(1/4)$ (Chudnovsky's theorem,
also a consequence of Nesterenko's theorem; see
\cite[p.~85]{Z}) implies algebraic independence of $\pi$ and $u$.
Thus distinct monomials in $u,\pi$ are linearly independent over
$\Q$. No combination of the printed terms can equal
\eqref{counter}. This proves the printed assertion false, while
\eqref{odd-result} supplies its proved correction.\qed

\section*{Task 2. Determine the exact nonzero support}
The theorem bounds possible monomials. Determine which coefficients
actually vanish for each denominator power, and whether all six
coefficients in the added odd band are nonzero for every $n\ge4$.
Seek formulas that also determine their signs.

\section*{Task 3. Treat every numerator degree}
Extend the divisor-interchange method to $x^{4p+1}$ for arbitrary
$p\ge0$, without a lower bound on $p$ in terms of $M$. Derive the
complete support restrictions and explain the exceptional vanishing
coefficients as uniform identities in modular forms.

\section*{Task 4. Resolve the first-moment conjecture}
Prove the two all-power formulas in Xu--Zhao Conjecture 5.8 using
the same method with numerator $x$. Determine whether the much
smaller number of derivatives gives closed recurrences for the
rational coefficients.

\section*{Task 5. Find an efficient recurrence in the power}
Develop a recurrence for the fifth-moment polynomials in $u,\pi$
as $M$ increases. Determine whether it can preserve the proved
support restrictions without first generating and cancelling the
larger residue expressions.

\begin{thebibliography}{4}
\bibitem{XZ} C.~Xu and J.~Zhao,
\emph{General Berndt-Type Integrals and Series Associated with Jacobi
Elliptic Functions},
\href{https://arxiv.org/abs/2401.01385v2}{arXiv:2401.01385v2} (2024).
\bibitem{Z} D.~Zagier, \emph{Elliptic Modular Forms and Their Applications},
in \emph{The 1-2-3 of Modular Forms}, Springer (2008), pp.~1--103,
Proposition 4 and p.~85;
\href{https://people.mpim-bonn.mpg.de/zagier/files/doi/10.1007/978-3-540-74119-0_1/fulltext.pdf}{author-hosted chapter}.
\bibitem{R} S.~Ramanujan, \emph{On certain arithmetical functions},
Trans. Cambridge Philos. Soc. \textbf{22} (1916), 159--184,
Section 7; \href{https://ramanujan.sirinudi.org/Volumes/published/ram18.html}{online transcription}.
\bibitem{DLMF} NIST Digital Library of Mathematical Functions,
\href{https://dlmf.nist.gov/20.7}{Section 20.7}, theta identities
and transformations.
\end{thebibliography}
\noindent\textit{Verification and status.}
The proof for arbitrary $M$ is the analytic and algebraic argument
above, not an inference from finite computation. Run
\texttt{verify\_evaluation.py} for the exact finite certificate and
separate numerical diagnostics. The full current source and dated
bounded claim checks accompany this note. The authoring agent
performed the proof audit; independent peer review has not occurred.
\end{document}
